\section{Обсуждение и ограничения}\label{sec:discussion}

Теорема и эксперименты дают согласованную картину. Постоянный шаг быстро
забывает начальное состояние, но создаёт смещение; PR-усреднение уменьшает
вариативность, не исправляя центр. RR использует зависимость двух ветвей
конструктивно: общая траектория сохраняет ведущий шум с коэффициентом
$2-1=1$, а комбинация шагов сокращает первый член разложения смещения.
Пуассоновское разложение затем превращает ведущую взвешенную сумму в
мартингал, а все негауссовские части собираются в один $L_p$-остаток.

Практически это означает следующее. Обе ветви нужно запускать на одной
траектории; выбор шага отвечает прежде всего за центр интервала, а оценка
долгосрочной дисперсии --- за его ширину. Эти две настройки следует
диагностировать раздельно.

У результата есть четыре ограничения.

\begin{itemize}
\item Теорема относится к фиксированным скалярным направлениям. Она не даёт
одновременной гарантии для всех координат и не является многомерной оценкой
Берри--Эссеена.
\item Финальное следствие использует треугольную схему
$\alpha=c/\sqrt n$, разогрев порядка $(\alpha a)^{-1}\log^2n$ и явное
условие $2\alpha\leq\alpha_{\mathrm{adm}}(p,q,\tmix)$. Оно не является
утверждением для фиксированного $\alpha$ без изменения центра.
\item Высокие моменты компонент разложения второго порядка импортируются из
\cite{Levin2026}. Поэтому допустимый шаг зависит от порядка момента $p$, от
$\tmix$ и от условия устойчивости случайных матричных произведений.
\item Близость OBM к оракулу установлена только эмпирически. Состоятельность
OBM для связанной RR-рекурсии, оптимальный размер блока и интервалы с
нормировкой, оценённой по данным, требуют отдельного анализа.
\end{itemize}

Эксперименты используют конечные цепи и умеренную зависимость. Проверка
функции распределения проведена для одной задачи и одного направления, а
популяционный тест --- для 100 задач и 100 траекторий на задачу; поэтому он
проверяет предсказанный механизм, но не заменяет стресс-тестов при медленном
смешивании.

\section*{Заключение}

Для PR-усреднённой после разогрева RR-оценки получена неасимптотическая
нормальная аппроксимация. В допустимом сбалансированном режиме
$\alpha=c/\sqrt n$, $p\asymp\log n$,
$n_0\asymp(\alpha a)^{-1}\log^2n$ и
$2\alpha\leq\alpha_{\mathrm{adm}}(p,q,\tmix)$ основной результат имеет вид
\[
 \dK\!\left(
 \frac{\sqrt n\,u^\top
  (\bar\theta_{n,n_0}^{\RR,\alpha}-\theta^\star)}
      {\sqrt{u^\top\Sigma_\infty u}},
 \mathcal N(0,1)\right)
 \lesssim
 \frac{C_{\mathrm{mix}}(\tmix)\operatorname{polylog}(n)}
      {n^{1/4}}.
\]
Доказательство выделяет ведущий пуассоновский мартингал, сравнивает его
детерминированную ковариационную прокси с $\Sigma_\infty$ и соединяет
опубликованные оценки стационарных компонент с локальной леммой переноса после
разогрева и леммой~\ref{lem:smoothing}.

В численных опытах одиночные ветви сохраняют заметно сдвинутый центр и
недопокрытие, тогда как RR достигает 94--95\%-го покрытия при практически той
же ширине. На длинных траекториях OBM близок к оракульной оценке дисперсии.
Таким образом, шаговая экстраполяция и оценивание долгосрочной дисперсии
являются дополняющими, но различными частями статистического вывода.

